\section{Experimental Validation}
\label{sec:experimental-validation}

The proposed Digitally Assisted Analog Front-End architecture was experimentally evaluated to verify that embedded digital assistance improves the effective operating conditions of the analog signal chain while preserving a conventional acquisition path during normal operation. The evaluation comprised foreground calibration on two independently assembled geophone acquisition boards under multiple gain configurations and quantitative residual DC-offset measurements.

\subsection{Prototype and Comparison Conditions}

The experimental validation comprised eight board--gain configurations: all five PGA settings on Board A together with the $\times16$, $\times32$, and $\times50$ settings on Board B. Each configuration was evaluated through ten paired restarts ($N=10$), representing repeated measurements of the same board and gain rather than independent devices.

A single fixed-reference baseline, $\mathbf{c}_{\mathrm{base}}$, selected empirically on Board B at $\times50$, was applied unchanged to all reported configurations. For each calibrated condition, the four analog stages were processed sequentially and the resulting DAC codes remained fixed during acquisition. Post-calibration EEPROM writes were disabled so that every restart began from the same baseline.

Board B at $\times1$ and $\times4$ were excluded because their baseline residuals already lay within the calibration deadbands, preventing corrective updates. A separate narrower-deadband diagnostic confirmed the expected behavior (Section~\ref{sec:pga-note}).

\subsection{Residual DC-Offset Measurement and Results}

The LP output was selected as the primary performance metric because it directly feeds the acquisition ADC. Measurements at the PGA, BP, and SUM outputs document the sequential propagation of the foreground calibration through the analog signal chain.

A four-channel Tektronix TDS~2004B recorded the mean output voltage of the PGA, BP, SUM, and LP stages over a 1-s window after settling. Probe grounds were connected to $V_{\mathrm{ref}}=V_{\mathrm{DDA}}/2$, so each channel measured $V_{\mathrm{stage}}-V_{\mathrm{ref}}$. Within each restart, the fixed-reference condition was always recorded before the calibrated condition, preserving paired measurements.

Table~\ref{tab:results} reports the mean and sample standard deviation of the normalized residual offset over ten paired restarts. Figure~\ref{fig:calibration_trace} shows representative calibration trajectories reconstructed by applying the 128-tap FIR used by the controller to oscilloscope records resampled at the controller rate.

\begin{figure}[t]
	\centering
	\includegraphics[width=0.95\linewidth]{imagenes/autocalibration_scope_sum.png}
	\caption{Representative foreground calibration trajectories: stage deviations relative to $V_{\mathrm{ref}}$ during calibration at $\times16$ gain (total differential gain $\times32$). Dashed light-gray traces are oscilloscope measurements; colored curves are resampled 128-tap-FIR reconstructions, not internal logs (Section~\ref{sec:pga-note}).}
	\label{fig:calibration_trace}
\end{figure}

Across the eight reported board--gain configurations, the proposed Digitally Assisted Analog Front-End consistently reduced the LP residual DC offset without requiring board-specific or gain-specific controller retuning. Relative to the shared fixed-reference baseline, the residual offset decreased from 2.9--45.8\%FS to 0.4--3.9\%FS. Individual calibrated observations ranged from 0.1--5.2\%FS (Table~\ref{tab:results}). A one-sided sign test on the paired LP measurements rejected the null hypothesis of no reduction in all eight configurations (10 of 10 paired restarts, exact $p\approx10^{-3}$).

\begin{table*}[!t]
	\caption{Residual DC offset (\%FS of the nominal 5-V ADC span), mean $\pm$ sample standard deviation over $N{=}10$ paired baseline-seeded restarts per reported board/gain cell. Within each cell, the upper and lower values correspond to the shared fixed-reference and calibrated conditions, respectively.}
	\label{tab:results}
	\centering
	\scriptsize
	\setlength{\tabcolsep}{3.2pt}
	\renewcommand{\arraystretch}{0.82}
	
	\resizebox{\textwidth}{!}{%
		\arrayrulecolor{gray!45}
		\begin{tabular}{l cc ;{2pt/2pt} cc ;{2pt/2pt} cccc}
			\arrayrulecolor{black}
			\toprule
			Stage
			& A, $\times1$
			& A, $\times4$
			& A, $\times16$
			& B, $\times16$
			& A, $\times32$
			& B, $\times32$
			& A, $\times50$
			& B, $\times50$ \\
			\midrule
			
			PGA
			& \shortstack{$3.5\pm0.2$\\[-1.5pt]$3.4\pm0.3$}
			& \shortstack{$2.4\pm0.1$\\[-1.5pt]$2.1\pm0.3$}
			& \shortstack{$2.4\pm0.1$\\[-1.5pt]$1.9\pm0.2$}
			& \shortstack{$2.1\pm0.1$\\[-1.5pt]$1.7\pm0.1$}
			& \shortstack{$1.0\pm0.2$\\[-1.5pt]$0.5\pm0.2$}
			& \shortstack{$1.0\pm0.2$\\[-1.5pt]$0.5\pm0.2$}
			& \shortstack{$1.2\pm0.3$\\[-1.5pt]$1.0\pm0.5$}
			& \shortstack{$2.5\pm0.1$\\[-1.5pt]$1.8\pm0.3$}
			\\
			
			\arrayrulecolor{gray!45}
			\hdashline[0.4pt/1.4pt]
			\arrayrulecolor{black}
			
			BP
			& \shortstack{$1.5\pm2.2$\\[-1.5pt]$1.2\pm0.2$}
			& \shortstack{$4.9\pm0.3$\\[-1.5pt]$1.0\pm0.2$}
			& \shortstack{$3.8\pm0.9$\\[-1.5pt]$1.5\pm1.5$}
			& \shortstack{$0.66\pm0.02$\\[-1.5pt]$0.66\pm0.02$}
			& \shortstack{$7.0\pm0.4$\\[-1.5pt]$1.0\pm0.5$}
			& \shortstack{$0.66\pm0.01$\\[-1.5pt]$0.66\pm0.01$}
			& \shortstack{$1.0\pm0.7$\\[-1.5pt]$0.9\pm0.5$}
			& \shortstack{$0.65\pm0.01$\\[-1.5pt]$0.65\pm0.02$}
			\\
			
			\arrayrulecolor{gray!45}
			\hdashline[0.4pt/1.4pt]
			\arrayrulecolor{black}
			
			SUM
			& \shortstack{$\mathbf{21.8\pm1.1}^{\dagger}$\\[-1.5pt]$0.8\pm0.2$}
			& \shortstack{$6.7\pm1.0$\\[-1.5pt]$1.0\pm0.3$}
			& \shortstack{$4.9\pm1.5$\\[-1.5pt]$1.4\pm1.5$}
			& \shortstack{$4.5\pm0.8$\\[-1.5pt]$0.6\pm0.5$}
			& \shortstack{$7.4\pm0.3$\\[-1.5pt]$0.6\pm0.2$}
			& \shortstack{$8.0\pm0.8$\\[-1.5pt]$1.2\pm0.3$}
			& \shortstack{$0.4\pm0.2$\\[-1.5pt]$\mathbf{0.8\pm0.6}^{\ddagger}$}
			& \shortstack{$0.9\pm0.6$\\[-1.5pt]$0.9\pm0.6$}
			\\
			
			\arrayrulecolor{gray!45}
			\hdashline[0.4pt/1.4pt]
			\arrayrulecolor{black}
			
			LP
			& \shortstack{$\mathbf{45.8\pm3.1}^{\dagger}$\\[-1.5pt]$3.2\pm0.9$}
			& \shortstack{$\mathbf{42.6\pm4.2}^{\dagger}$\\[-1.5pt]$3.9\pm0.8$}
			& \shortstack{$\mathbf{29.6\pm6.8}^{\dagger}$\\[-1.5pt]$2.3\pm1.4$}
			& \shortstack{$3.0\pm0.1$\\[-1.5pt]$0.7\pm0.1$}
			& \shortstack{$\mathbf{12.4\pm1.3}^{\dagger}$\\[-1.5pt]$1.7\pm0.7$}
			& \shortstack{$2.9\pm0.1$\\[-1.5pt]$0.4\pm0.2$}
			& \shortstack{$4.1\pm2.1$\\[-1.5pt]$1.2\pm0.4$}
			& \shortstack{$2.9\pm0.1$\\[-1.5pt]$0.7\pm0.2$}
			\\
			
			\bottomrule
		\end{tabular}%
		\arrayrulecolor{black}
	}
	
	\vspace{1pt}
	\footnotesize
	Lower is better.
	$^{\dagger}$Fixed reference $\geq10\%$FS.
	$^{\ddagger}$The marked calibrated mean is larger than its baseline;
	equality after rounding is not marked
	(Section~\ref{sec:pga-note}).
\end{table*}

